\section{Termination by descending boundary coefficients}
\label{sec:termination}

We now combine the two counts.  At each positive boundary coefficient,
the witness count first excludes flipped surfaces, and the cycle count
then excludes flipping surfaces.  The same witness count at threshold
two removes all remaining flipped surfaces, after which the ranks of
surface classes on the fourfolds decrease at every step.

Write $L_i$ and $L_i^+$ for the exceptional loci of step $i$, so that
flipping surfaces lie in $L_i$ and flipped surfaces lie in $L_i^+$.
Proposition~\ref{prop:comparison} gives
\begin{equation}\label{eq:final-dimension}
 \dim L_i\leq2,\qquad \dim L_i^+\leq2,\qquad
 \dim L_i+\dim L_i^+\geq3.
\end{equation}
Smallness identifies the finitely many positive boundary
components and preserves their coefficients.

\begin{lemma}[Persistence of a surface centre]\label{lem:persistence}
Let $T\dashrightarrow T^+$ be a step of a sequence as in
Theorem~\ref{thm:termination}, and let $E$ be an exceptional place
whose centre on $T$ is a surface $S$.  If $S$ is not contained in the
source exceptional locus, its target centre is the strict transform
of $S$, a surface meeting the common isomorphic open. If $S$
lies in a boundary component, its transform lies in the
corresponding component.
\end{lemma}

\begin{proof}
Represent the place on a common higher model.  Its image onto $S$ is
surjective by the definition of the centre and properness.  The inverse
image of the dense open portion of $S$ on which the step is an
isomorphism is nonempty and dense in $E$.  There the two centre maps agree.
The closure of this portion's target image is both the strict
transform of $S$ and the whole target centre. The assertions
follow.
\end{proof}

The following count isolates the argument used at each coefficient,
including the final value $b=0$.  Write
$a_i(E)=a(E;X_i,B_i+\bM)$ and fix the integer $m$ of
Proposition~\ref{prop:witnesses}.

\begin{lemma}[Counting flipped surfaces]
\label{lem:coefficient-witness-count}
Consider a tail beginning with $X_N$ of a sequence as in
Theorem~\ref{thm:termination}.  Let $b$ be either zero or a positive
boundary coefficient.  Suppose no flipping surface at any step of the
tail lies in a boundary component of coefficient greater than $b$.
If $b>0$, only finitely many steps of the tail have a flipped surface
contained in a coefficient-$b$ boundary component.  If $b=0$, only
finitely many steps of the tail have any flipped surface.
\end{lemma}

\begin{proof}
At every step $i$ under consideration, choose one such flipped surface
and its witness $E_i$ from Proposition~\ref{prop:witnesses}.  The place
is exceptional over $X_N$, and discrepancy monotonicity gives
\begin{equation}\label{eq:tail-witness-bound}
 a_N(E_i)\leq a_i(E_i)<a_{i+1}(E_i)\leq2-b.
\end{equation}
For $b=0$ this is the witness bound for an arbitrary flipped surface.

Apply Lemma~\ref{lem:low-places} on the fixed model $X_N$.  If
$E_i\notin\mathcal F_{X_N}$, its centre there is a surface contained
in a positive boundary component of coefficient $d$, and
\[
 2-d\leq a_N(E_i)<2-b.
\]
Thus $d>b$.  By the hypothesis on flipping surfaces, this surface
centre is not contained in $L_N$.  Lemma~\ref{lem:persistence} carries
it to a surface in the corresponding coefficient-$d$ component on
$X_{N+1}$.  At each successive step up to and including step $i$, the
same hypothesis and lemma apply.  Consequently the centre on $X_{i+1}$
meets $X_{i+1}\setminus L_i^+$.  This contradicts its definition as
the chosen flipped surface, which is contained in $L_i^+$.

Every chosen witness therefore belongs to the fixed finite set
$\mathcal F_{X_N}$.  Lemma~\ref{lem:finite-use} permits each place to
be used at most $m$ times, since its discrepancies on witness targets
are strictly increasing members of $(1,2]\cap m^{-1}\Z$.  There are
at most $m\,\#\mathcal F_{X_N}$ steps under consideration.
\end{proof}

\begin{proof}[Proof of Theorem~\ref{thm:termination}]
Suppose the sequence is infinite.  Proposition~\ref{prop:comparison}
preserves generalized canonicality: every exceptional place has
log discrepancy at least one on every model.  The strict bound greater
than one used in the witness count holds on the target of its witness
step by strict discrepancy increase.  No terminality of the whole
model is used.

Order the distinct positive boundary coefficients as
$b_1>\cdots>b_s>0$.  Inductively pass to tails on which no flipping
or flipped surface lies in a component of any coefficient already
treated.  Start with the full sequence.  For the next coefficient
$b=b_j$, the induction hypothesis satisfies the assumption of
Lemma~\ref{lem:coefficient-witness-count}.  Passing beyond the finitely
many steps it counts leaves a tail with no flipped surface in any
coefficient-$b$ component.

Index those components on the initial model by a finite set $I_b$,
and write $H_{i,\alpha}$ for their strict transforms on $X_i$,
$\alpha\in I_b$.  Smallness preserves this indexing.  Each target
component now satisfies
\[
 \dim(H_{i+1,\alpha}\cap L_i^+)\leq1.
\]
Otherwise this intersection would contain a flipped surface.
The boundary-component application of Lemma~\ref{lem:cycle-loss}
therefore gives
\[
 c_2(H_{i,\alpha})\geq c_2(H_{i+1,\alpha}),
\]
with strict inequality whenever $H_{i,\alpha}$ contains a flipping
surface.  The hypotheses were checked in Section~\ref{sec:cycles}
for the reduced components, without normalizing them.  In particular,
K\"ahler positivity in $X_i$ makes every removed surface class nonzero
in $H_4(H_{i,\alpha},\R)$.

It follows that the finite nonnegative integer
\[
 r_b(i):=\sum_{\alpha\in I_b}c_2(H_{i,\alpha})
\]
does not increase and decreases strictly at each step with a flipping
surface in a coefficient-$b$ component.  Only finitely many such
steps occur.  Pass beyond them to finish this stage.

After all positive coefficients have been treated, the hypothesis of
Lemma~\ref{lem:coefficient-witness-count} holds with $b=0$.  If
$B_0=0$, the positive-coefficient induction is empty and this hypothesis
already holds on the original sequence.  Pass beyond the finitely many
steps having a flipped surface.  On the resulting tail,
$\dim L_i^+\leq1$, whereas \eqref{eq:final-dimension} forces $L_i$ to
contain a surface.  Lemma~\ref{lem:cycle-loss}, applied to the
fourfolds themselves, gives
\[
 c_2(X_i)>c_2(X_{i+1})
\]
at every remaining step.  These are finite nonnegative integers, so
they cannot decrease strictly along an infinite tail.  This proves
termination.
\end{proof}
